\chapter{Introduction}

\section{Overview}

This document highlights the \textbf{technical aspects} of the penetration testing activities performed by the student \textbf{Mariano Aponte} to assess the security posture of the \textbf{"Corrosion: 1"} machine from \textit{VulnHub}.
\\

This project was carried out within the scope of the \textbf{2025--2026 Penetration Testing and Ethical Hacking} examination for the Master's Degree in Computer Science, Cybersecurity curriculum, taught by \textbf{Prof. Arcangelo Castiglione}.

\section{Methodologies}

To assess the security posture of the target system, a penetration testing methodology had to be selected. In this project, the chosen methodology is the \textbf{\textit{Penetration Testing Execution Standard (PTES)}} \cite{ptes}.
\\

Throughout this document, the penetration testing process follows the PTES methodology as meticulously as possible. Nevertheless, the phases defined by the standard are not intended to provide a comprehensive set of instructions for conducting a penetration test. Instead, they serve as guidelines that direct the assessment process, as stated in the PTES technical guidelines.
\\

\subsection{PTES Phases}

The \textit{Penetration Testing Execution Standard} defines the following seven phases for conducting a penetration test:

\begin{itemize}
    \item Pre-Engagement Interactions;
    \item Intelligence Gathering;
    \item Threat Modeling;
    \item Vulnerability Analysis;
    \item Exploitation;
    \item Post Exploitation;
    \item Reporting.
\end{itemize}

This report is organized according to the PTES methodology and focuses on the activities performed during the \textbf{Intelligence Gathering}, \textbf{Vulnerability Analysis}, \textbf{Exploitation}, and \textbf{Post Exploitation} phases. 
\\

Each phase is described in detail in the following chapters.

\section{Pentesting Tools}

\subsection{Environment}

To evaluate the security posture of the target asset, \textbf{"Corrosion: 1"} \cite{vulnhubCorrosion}, a dedicated virtual laboratory was set up.
\\

The laboratory was built using \texttt{Oracle VM VirtualBox 7.2.8 r173730} \cite{virtualbox}. A \texttt{NAT Network} was created, and all virtual machines deployed within VirtualBox were connected to this isolated network.
\\

The laboratory network was configured with the address space \texttt{10.0.2.0/24}, with \texttt{DHCP} enabled to automatically assign IP addresses to the virtual machines.
\\

The target machine, \textbf{"Corrosion: 1"}, was downloaded from the \textit{VulnHub} platform and deployed as a virtual machine using VirtualBox.
\\

The attacking machine was \texttt{Kali Linux 2026.2} \cite{kali}, downloaded from the official \textit{kali.org} website and deployed as a virtual machine within the same private network as the target. All penetration testing activities described in this report were performed from this system.

\subsection{Tools}

The following tools were used throughout the penetration testing activity:

\begin{itemize}
	\item Oracle VirtualBox \cite{virtualbox};
	\item Kali Linux \cite{kali};
	\item \texttt{nmap} \cite{nmap};
	\item \texttt{DirBuster} \cite{dirbuster};
	\item \texttt{Gobuster} \cite{gobuster};
	\item \texttt{wfuzz} \cite{wfuzz};
	\item \texttt{Nessus} \cite{nessus};
	\item \texttt{Nikto} \cite{nikto};
	\item \texttt{WafW00f} \cite{wafwoof};
	\item \texttt{Exploit Database} \cite{exploitdb};
	\item \texttt{Metasploit} \cite{metasploitMetasploitPenetration};
	\item \texttt{ssh} and \texttt{sftp} \cite{ssh};
	\item \texttt{curl} \cite{curl};
	\item \texttt{zip2john} and \texttt{john} \cite{john};
	\item \texttt{C} \cite{c}, \texttt{PHP} \cite{php} and \texttt{Python} \cite{python} programming;
	\item \texttt{netcat} \cite{netcat}.
\end{itemize}



